\chapter{Complexity — Prime Find}
%%% generated by the blueprint pipeline from TCSlib.Complexity.ClassP.BinaryArith.PrimeFind
\section{Overview}

Declarations of the Lean module \texttt{TCSlib.Complexity.ClassP.BinaryArith.PrimeFind}.

\section{Declarations}

\begin{definition}[Lovász bit bound]
\lean{Complexity.BinArith.lovaszBits}
\label{Complexity.BinArith.lovaszBits}
\leanok
For a natural number $n$, the Lovász bit bound is the least natural number $k$ such that
$2^k
\ge 2n$.
\end{definition}

\begin{theorem}[Power bound for the Lovász bit count]
\lean{Complexity.BinArith.two_mul_le_two_pow_lovaszBits}
\label{Complexity.BinArith.two_mul_le_two_pow_lovaszBits}
\leanok
\uses{Complexity.BinArith.lovaszBits}
\difficulty{1}
For every natural number $n$, if $L=\lceil\log_2(2n)\rceil$ (with the natural-number
ceiling convention), then $2n \le 2^L$.
\end{theorem}

\begin{theorem}[Bound on the Lovász bit count]
\lean{Complexity.BinArith.lovaszBits_le}
\label{Complexity.BinArith.lovaszBits_le}
\leanok
\uses{Complexity.BinArith.lovaszBits}
\difficulty{3}
For every natural number $n$, the ceiling of the base-$2$ logarithm of $2n$ is at most
$n+1$.
\end{theorem}

\begin{theorem}[Strict power bound for the Lovász bit count]
\lean{Complexity.BinArith.two_pow_lovaszBits_lt}
\label{Complexity.BinArith.two_pow_lovaszBits_lt}
\leanok
\uses{Complexity.BinArith.lovaszBits}
\difficulty{4}
For every nonzero natural number $n$, if $L=\lceil\log_2(2n)\rceil$ (with the
natural-number ceiling convention), then $2^L < 4n$.
\end{theorem}

\begin{theorem}[Prime in a dyadic interval]
\lean{Complexity.BinArith.exists_prime_lovaszBits}
\label{Complexity.BinArith.exists_prime_lovaszBits}
\leanok
\uses{Complexity.BinArith.lovaszBits}
\difficulty{4}
For every natural number $n$, there is a prime number $p$ strictly greater than
$2^{\ell(n)}$ and at most $2^{\ell(n)+1}$, where $\ell(n)$ is the least natural number
$k$ such that $2^k\ge 2n$.
\end{theorem}

\begin{definition}[Least prime above the binary size bound]
\lean{Complexity.BinArith.lovaszPrime}
\label{Complexity.BinArith.lovaszPrime}
\leanok
For each natural number $n$, define $p(n)$ to be the least prime strictly greater than
$2^{\lceil\log_2(2n)\rceil}+1$, where $\lceil\log_2(2n)\rceil$ is the least natural
number $k$ such that $2n \le 2^k$.
\end{definition}

\begin{theorem}[Primality of the chosen number]
\lean{Complexity.BinArith.lovaszPrime_prime}
\label{Complexity.BinArith.lovaszPrime_prime}
\leanok
\uses{Complexity.BinArith.lovaszPrime}
\difficulty{1}
For every natural number $n$, the number chosen as the least prime greater than
$2^{\lceil\log_2(2n)\rceil}+1$ is prime.
\end{theorem}

\begin{theorem}[Power bound for the least prime]
\lean{Complexity.BinArith.two_pow_lt_lovaszPrime}
\label{Complexity.BinArith.two_pow_lt_lovaszPrime}
\leanok
\uses{Complexity.BinArith.lovaszPrime}
\difficulty{1}
For every natural number $n$, let $p_n$ be the least prime number greater than or equal
to $2^{\lceil\log_2(2n)\rceil}+1$, with the ceiling taken in the natural-number
convention. Then $2^{\lceil\log_2(2n)\rceil} < p_n$.
\end{theorem}

\begin{theorem}[Primes below the Lovász prime]
\lean{Complexity.BinArith.not_prime_of_lt_lovaszPrime}
\label{Complexity.BinArith.not_prime_of_lt_lovaszPrime}
\leanok
\uses{Complexity.BinArith.lovaszPrime}
\difficulty{1}
For each natural number $n$, let the Lovász prime be the least prime strictly greater
than $2^{\lceil\log_2(2n)\rceil}+1$. If a natural number $q$ satisfies
$2^{\lceil\log_2(2n)\rceil}<q$ and $q$ is less than the Lovász prime, then $q$ is not
prime.
\end{theorem}

\begin{theorem}[Upper bound for the chosen prime]
\lean{Complexity.BinArith.lovaszPrime_le}
\label{Complexity.BinArith.lovaszPrime_le}
\leanok
\uses{Complexity.BinArith.lovaszPrime}
\difficulty{4}
For each natural number $n$, let $k$ be the least natural number such that $2n\le2^k$ (with $k=0$ when $n=0$). The least prime strictly greater than $2^k$ is at most $2^{k+1}$.
\end{theorem}

\begin{theorem}[Upper bound for a power of two]
\lean{Complexity.BinArith.two_pow_clog_le}
\label{Complexity.BinArith.two_pow_clog_le}
\leanok
\difficulty{5}
For every natural number $n$, $2^{\lceil\log_2(2n)\rceil} \le 4n+1$, where the ceiling
is taken in the natural-number convention.
\end{theorem}

\begin{theorem}[Linear bound for the chosen prime]
\lean{Complexity.BinArith.lovaszPrime_le_linear}
\label{Complexity.BinArith.lovaszPrime_le_linear}
\leanok
\uses{Complexity.BinArith.lovaszPrime, Complexity.BinArith.lovaszPrime_le, Complexity.BinArith.two_pow_clog_le}
\difficulty{3}
For every natural number $n$, the prime chosen as the least prime greater than
$2^{\lceil\log_2(2n)\rceil}+1$ is at most $8n+2$.
\end{theorem}

\begin{definition}[Counting step]
\lean{Complexity.BinArith.countStep}
\label{Complexity.BinArith.countStep}
\leanok
\uses{Complexity.BinArith.add, Complexity.isNilB, Complexity.pairFstD, Complexity.pairSndD, Turing.pairEncode}
For a list state $S$, if its second component is empty, the step returns $S$. Otherwise,
writing its defaulted pair decomposition as $(A,B)$, the step encodes the pair
consisting of the binary sum of the encoding of $(A,[	ext{true}])$ and the list obtained
by dropping the first element of $B$.
\end{definition}

\begin{definition}[Count loop]
\lean{Complexity.BinArith.countLoop}
\label{Complexity.BinArith.countLoop}
\leanok
\uses{Complexity.BinArith.countStep, Turing.pairEncode}
For a finite sequence $z$ of Boolean values, the count loop is the result of applying
the counting step $(	ext{length of }z+1)$ times to the encoding of the pair consisting
of the empty sequence and $z$.
\end{definition}

\begin{definition}[Counting invariant]
\lean{Complexity.BinArith.CountInv}
\label{Complexity.BinArith.CountInv}
\leanok
\uses{Complexity.LogProg.bitsVal, Turing.pairEncode}
For lists $z$ and $S$, the counting invariant holds if there are lists $a$ and $r$ such
that $S$ encodes $(a,r)$, $a$ is the canonical binary representation of its value, and
the binary value of $a$ plus the length of $r$ equals the length of $z$.
\end{definition}

\begin{theorem}[Invariant preservation and step decrease]
\lean{Complexity.BinArith.countInv_step}
\label{Complexity.BinArith.countInv_step}
\leanok
\uses{Complexity.BinArith.CountInv, Complexity.BinArith.add, Complexity.BinArith.countStep, Complexity.LogProg.bitsVal, Complexity.LogProg.bitsVal_bits, Complexity.isNilB, Complexity.pairSndD, Complexity.pairSndD_pairEncode, Turing.pairEncode}
\difficulty{5}
Let $z$ and $S$ be bit strings, and suppose there are bit strings $a,r$ such that $S$ is
the pair encoding of $a$ and $r$, $a$ is the binary bit list of its own value, and the
value of $a$ plus the length of $r$ equals the length of $z$. After one counting step,
these properties still hold; if the decoded second component of $S$ has positive length,
its length strictly decreases, and if it has length zero, the step leaves $S$ unchanged.
\end{theorem}

\begin{theorem}[Counting loop output]
\lean{Complexity.BinArith.countLoop_eq}
\label{Complexity.BinArith.countLoop_eq}
\leanok
\uses{Complexity.BinArith.CountInv, Complexity.BinArith.countInv_step, Complexity.BinArith.countLoop, Complexity.BinArith.countStep, Complexity.BinArith.iterate_invariant, Complexity.BinArith.iterate_measure_eq_zero, Complexity.LogProg.bitsVal, Complexity.pairSndD, Complexity.pairSndD_pairEncode, Turing.pairEncode}
\difficulty{5}
For every bit string $z$, iterating the counting step $|z|+1$ times, starting from the
pair encoding of the empty word and $z$, gives the pair encoding of the binary bit list
of $|z|$ and the empty word.
\end{theorem}

\begin{theorem}[Polynomial-time counting loop]
\lean{Complexity.BinArith.polyTimeComputable_countLoop}
\label{Complexity.BinArith.polyTimeComputable_countLoop}
\leanok
\uses{Complexity.BinArith.CountInv, Complexity.BinArith.countInv_step, Complexity.BinArith.countLoop, Complexity.BinArith.countStep, Complexity.BinArith.iterate_invariant, Complexity.BinArith.polyTimeComputable_add, Complexity.BinArith.polyTimeComputable_iterate, Complexity.LogProg.bitsVal, Complexity.LogProg.length_bits_le, Complexity.PolyTimeComputable, Complexity.PolyTimeComputable.comp, Complexity.PolyTimeComputable.pairEncode, Complexity.polyTimeComputable_const, Complexity.polyTimeComputable_id, Complexity.polyTimeComputable_isNil, Complexity.polyTimeComputable_ite, Complexity.polyTimeComputable_pairFstD, Complexity.polyTimeComputable_pairSndD, Complexity.polyTimeComputable_tail, Turing.length_pairEncode, Turing.pairEncode}
\difficulty{5}
The function that starts with the pair encoding of the empty word and an input word,
then iterates the following update one more than the input length, is polynomial-time
computable: if the second projection is empty, leave the pair unchanged; otherwise,
replace it by the pair encoding of the sum of the decoded first projection and the
one-bit word $[1]$, together with the second projection after its first bit is removed.
\end{theorem}

\begin{theorem}[Value of a terminal one bit]
\lean{Complexity.BinArith.bitsVal_replicate_false_one}
\label{Complexity.BinArith.bitsVal_replicate_false_one}
\leanok
\uses{Complexity.BinArith.bitsVal_append_singleton, Complexity.BinArith.bitsVal_eq_zero_iff, Complexity.LogProg.bitsVal}
\difficulty{1}
For every natural number $j$, the binary value of the word consisting of $j$ false bits
followed by one true bit is $2^j$.
\end{theorem}

\begin{definition}[Conditional doubling step]
\lean{Complexity.BinArith.dblStep}
\label{Complexity.BinArith.dblStep}
\leanok
\uses{Complexity.BinArith.lt, Complexity.pairFstD, Complexity.pairSndD, Turing.pairEncode}
For a Boolean list $S$, define its next state to be the pair encoding of the list
obtained by prefixing `false` to the first component of the default first-pair
decomposition of $S$, together with the second component of that decomposition, if $S$
satisfies the predicate `lt`; otherwise, leave $S$ unchanged.
\end{definition}

\begin{definition}[Iterated doubling loop]
\lean{Complexity.BinArith.dblLoop}
\label{Complexity.BinArith.dblLoop}
\leanok
\uses{Complexity.BinArith.countLoop, Complexity.BinArith.dblStep, Complexity.pairFstD, Turing.pairEncode}
For a Boolean list $z$, define the output by iterating the conditional doubling step
$2(|z|+1)^1$ times, starting from the pair encoding of $[\mathrm{true}]$ and the list
formed by prefixing `false` to the first component of the default first-pair
decomposition of the output of the counting loop on $z$.
\end{definition}

\begin{definition}[Doubling invariant]
\lean{Complexity.BinArith.DblInv}
\label{Complexity.BinArith.DblInv}
\leanok
\uses{Turing.pairEncode}
For a natural number $n$ and a Boolean list $S$, the doubling invariant holds when there
is a natural number $j$ such that $S$ is the pair encoding of the list consisting of $j$
copies of `false` followed by `true`, and the list formed by prefixing `false` to the
binary-bit list of $n$, with $j\leq \lceil\log_2(2n)\rceil$.
\end{definition}

\begin{theorem}[Doubling invariant under one step]
\lean{Complexity.BinArith.dblInv_step}
\label{Complexity.BinArith.dblInv_step}
\leanok
\uses{Complexity.BinArith.DblInv, Complexity.BinArith.bitsVal_replicate_false_one, Complexity.BinArith.dblStep, Complexity.BinArith.lt, Complexity.LogProg.bitsVal, Complexity.LogProg.bitsVal_bits, Complexity.pairFstD, Complexity.pairFstD_pairEncode, Complexity.pairSndD_pairEncode, Turing.pairEncode}
\difficulty{6}
Let $n$ be a natural number and $S$ a bit string. Suppose some $j$ satisfies that $S$
encodes the pair consisting of $j$ zero bits followed by one, and a zero followed by the
binary bit list of $n$, with $j\leq \lceil\log_2(2n)\rceil$. One doubling step preserves
this property. If the potential $\lceil\log_2(2n)\rceil-(|a|-1)$ is positive, where $a$
is the first decoded component of $S$, then the same potential for the state after the
step is strictly smaller; if it is zero, the step leaves $S$ unchanged.
\end{theorem}

\begin{theorem}[First component after doubling loop]
\lean{Complexity.BinArith.dblLoop_spec}
\label{Complexity.BinArith.dblLoop_spec}
\leanok
\uses{Complexity.BinArith.DblInv, Complexity.BinArith.countLoop, Complexity.BinArith.countLoop_eq, Complexity.BinArith.dblInv_step, Complexity.BinArith.dblLoop, Complexity.BinArith.dblStep, Complexity.BinArith.iterate_invariant, Complexity.BinArith.iterate_measure_eq_zero, Complexity.BinArith.lovaszBits_le, Complexity.pairFstD, Complexity.pairFstD_pairEncode, Turing.pairEncode}
\difficulty{4}
For every bit string $z$, the first decoded component after the doubling loop is the
word consisting of $\lceil\log_2(2|z|)\rceil$ zero bits followed by one.
\end{theorem}

\begin{theorem}[Polynomial-time computability of the doubling loop]
\lean{Complexity.BinArith.polyTimeComputable_dblLoop}
\label{Complexity.BinArith.polyTimeComputable_dblLoop}
\leanok
\uses{Complexity.BinArith.DblInv, Complexity.BinArith.countLoop, Complexity.BinArith.countLoop_eq, Complexity.BinArith.dblInv_step, Complexity.BinArith.dblLoop, Complexity.BinArith.dblStep, Complexity.BinArith.iterate_invariant, Complexity.BinArith.lovaszBits_le, Complexity.BinArith.polyTimeComputable_cons, Complexity.BinArith.polyTimeComputable_countLoop, Complexity.BinArith.polyTimeComputable_iterate, Complexity.BinArith.polyTimeComputable_lt, Complexity.LogProg.length_bits_le, Complexity.PolyTimeComputable, Complexity.PolyTimeComputable.comp, Complexity.PolyTimeComputable.pairEncode, Complexity.pairFstD, Complexity.polyTimeComputable_const, Complexity.polyTimeComputable_id, Complexity.polyTimeComputable_ite, Complexity.polyTimeComputable_pairFstD, Complexity.polyTimeComputable_pairSndD, Turing.length_pairEncode, Turing.pairEncode}
\difficulty{4}
The function on bit strings obtained by applying the doubling step $2(|z|+1)$ times to
the encoded pair whose first component is $1$ and whose second component is the
concatenation of $0$ with the first component of the result of applying the counting
step $|z|+1$ times to the encoded pair of the empty string and $z$ is polynomial-time
computable.
\end{theorem}

\begin{definition}[Binary arithmetic search step]
\lean{Complexity.BinArith.searchStep}
\label{Complexity.BinArith.searchStep}
\leanok
\uses{Complexity.BinArith.add, Complexity.BinArith.beq, Complexity.BinArith.mod, Complexity.LogProg.bitsVal, Complexity.pairFstD, Complexity.pairSndD, Turing.pairEncode}
For a Boolean list $S$, define a new list by the following cases. If the default-false
head of the second component of the second default pair decomposition of $S$ is true,
the new list is $S$. Otherwise, if the Boolean test on the pair encoding of the first
component of the second default pair decomposition of $S$ and the first component of $S$
holds, the new list is the pair encoding of the first component of $S$ with the pair
encoding of that former component and $[\mathrm{true}]$. If neither condition holds and
the bit value of the specified modular expression is zero, the new list is the pair
encoding of the incremented pair encoding of the first component of $S$ with
$[\mathrm{true}]$, and the pair encoding of $[\mathrm{false},\mathrm{true}]$ with
$[\mathrm{false}]$. In the remaining case, it is the pair encoding of the first
component of $S$ with the pair encoding of the incremented pair encoding of the first
component of the second default pair decomposition of $S$ with $[\mathrm{true}]$, and
$[\mathrm{false}]$.
\end{definition}

\begin{definition}[Search invariant]
\lean{Complexity.BinArith.SearchInv}
\label{Complexity.BinArith.SearchInv}
\leanok
\uses{Complexity.BinArith.lovaszPrime, Turing.pairEncode}
For a natural number $n$ and a Boolean list $S$, the search invariant holds exactly when
there are natural numbers $c,d$ and a Boolean $b$ such that $S$ is the pair encoding of
the binary expansion of $c$ and the pair encoding of the binary expansion of $d$ with
the one-element list containing $b$, and
$2^{\lceil\log_2(2n)\rceil}<c\leq \operatorname{lovaszPrime}(n)$, $2\leq d\leq c$, no
natural number $m$ with $2\leq m<d$ divides $c$, and $b=\mathrm{true}$ implies
$c=\operatorname{lovaszPrime}(n)$.
\end{definition}

\begin{definition}[Search measure]
\lean{Complexity.BinArith.searchMeasure}
\label{Complexity.BinArith.searchMeasure}
\leanok
\uses{Complexity.BinArith.lovaszPrime, Complexity.LogProg.bitsVal, Complexity.pairFstD, Complexity.pairSndD}
For a natural number $n$ and a Boolean list $S$, let $a$ be the first component of the
default pair decomposition of $S$, let $b$ be the first component of the default pair
decomposition of the second component of $S$, and let $	au$ be the first entry of the
second component of that second component, taking $	au$ to be false if the component is
empty. Define the search measure to be $0$ when $	au$ is true; otherwise define it as
$(P(n)-v(a))(P(n)+1)+(v(a)-v(b))+1$, where $P(n)$ is the prime-valued function used in
the definition and $v$ is the binary value of a Boolean list.
\end{definition}

\begin{theorem}[Preservation and Progress of the Search Invariant]
\lean{Complexity.BinArith.searchInv_step}
\label{Complexity.BinArith.searchInv_step}
\leanok
\uses{Complexity.BinArith.SearchInv, Complexity.BinArith.add, Complexity.BinArith.beq, Complexity.BinArith.lovaszPrime, Complexity.BinArith.lovaszPrime_prime, Complexity.BinArith.mod, Complexity.BinArith.not_prime_of_lt_lovaszPrime, Complexity.BinArith.searchMeasure, Complexity.BinArith.searchStep, Complexity.LogProg.bitsVal, Complexity.LogProg.bitsVal_bits, Complexity.pairFstD_pairEncode, Complexity.pairSndD_pairEncode, Turing.pairEncode}
\difficulty{6}
Let $n$ be a natural number and $S$ a bit string. Suppose there are natural numbers
$c,d$ and a bit $b$ such that
$S=\langle\operatorname{bits}(c),\langle\operatorname{bits}(d),[b]\rangle\rangle$, where
$2^{\lceil\log_2(2n)\rceil}<c\leq q$, $2\leq d\leq c$, no integer $m$ with $2\leq m<d$
divides $c$, and $b$ can be true only when $c=q$; here $q$ is the least prime greater
than $2^{\lceil\log_2(2n)\rceil}$. Then one search step preserves these conditions. It
strictly decreases the specified search measure whenever that measure is positive, and
leaves $S$ unchanged whenever that measure is zero.
\end{theorem}

\begin{definition}[Search initialization]
\lean{Complexity.BinArith.searchInit}
\label{Complexity.BinArith.searchInit}
\leanok
\uses{Complexity.BinArith.add, Complexity.BinArith.dblLoop, Complexity.pairFstD, Turing.pairEncode}
For a Boolean list $z$, define the initial search state as the pair encoding of the
increment of the first component of the output of the doubling loop on $z$, and the pair
encoding of $[\mathrm{false},\mathrm{true}]$ with $[\mathrm{false}]$.
\end{definition}

\begin{theorem}[Initial Search Invariant]
\lean{Complexity.BinArith.searchInv_init}
\label{Complexity.BinArith.searchInv_init}
\leanok
\uses{Complexity.BinArith.SearchInv, Complexity.BinArith.add, Complexity.BinArith.bitsVal_replicate_false_one, Complexity.BinArith.dblLoop_spec, Complexity.BinArith.searchInit, Complexity.BinArith.two_pow_lt_lovaszPrime, Complexity.LogProg.bitsVal}
\difficulty{4}
For every bit string $z$, there exist natural numbers $c,d$ and a bit $b$ such that the
initial state for input $z$ encodes $c,d,b$, with $2^{\lceil\log_2(2|z|)\rceil}<c\leq
q$, $2\leq d\leq c$, and no integer $m$ with $2\leq m<d$ divides $c$. Here $q$ is the
least prime greater than $2^{\lceil\log_2(2|z|)\rceil}$. The bit $b$ can be true only
when $c=q$.
\end{theorem}

\begin{definition}[Search iteration loop]
\lean{Complexity.BinArith.searchLoop}
\label{Complexity.BinArith.searchLoop}
\leanok
\uses{Complexity.BinArith.searchInit, Complexity.BinArith.searchStep}
For a Boolean list $z$, define the output by iterating the search state transition
$64(|z|+1)^2$ times, starting from the initial search state on $z$.
\end{definition}

\begin{theorem}[Polynomial-Time Computability of a Prime Encoding]
\lean{Complexity.BinArith.polyTimeComputable_lovaszPrime}
\label{Complexity.BinArith.polyTimeComputable_lovaszPrime}
\leanok
\uses{Complexity.BinArith.SearchInv, Complexity.BinArith.iterate_invariant, Complexity.BinArith.iterate_measure_eq_zero, Complexity.BinArith.lovaszPrime, Complexity.BinArith.lovaszPrime_le_linear, Complexity.BinArith.polyTimeComputable_add, Complexity.BinArith.polyTimeComputable_beq, Complexity.BinArith.polyTimeComputable_dblLoop, Complexity.BinArith.polyTimeComputable_isZero, Complexity.BinArith.polyTimeComputable_iterate, Complexity.BinArith.polyTimeComputable_mod, Complexity.BinArith.polyTimeComputable_slot2, Complexity.BinArith.polyTimeComputable_slot3, Complexity.BinArith.searchInit, Complexity.BinArith.searchInv_init, Complexity.BinArith.searchInv_step, Complexity.BinArith.searchLoop, Complexity.BinArith.searchMeasure, Complexity.BinArith.searchStep, Complexity.LogProg.bitsVal_bits, Complexity.LogProg.length_bits_le, Complexity.LogProg.length_bits_mono, Complexity.PolyTimeComputable, Complexity.PolyTimeComputable.comp, Complexity.PolyTimeComputable.pairEncode, Complexity.pairFstD_pairEncode, Complexity.pairSndD_pairEncode, Complexity.polyTimeComputable_const, Complexity.polyTimeComputable_headD, Complexity.polyTimeComputable_id, Complexity.polyTimeComputable_ite, Complexity.polyTimeComputable_pairFstD, Turing.length_pairEncode}
\difficulty{6}
For a bit string $z$, let $p(z)$ be the least prime at least
$2^{\lceil\log_2(2|z|)\rceil}+1$, and write its binary expansion as a bit string. The
function sending $z$ to this expansion is computable in polynomial time.
\end{theorem}
